% Compile with: tectonic "C++_Struct與Class_完整參考筆記.tex"
\documentclass[11pt,a4paper]{article}

\usepackage[a4paper,top=22mm,bottom=22mm,left=23mm,right=23mm,headheight=22pt]{geometry}
\usepackage{fontspec}
\usepackage{xeCJK}
\usepackage{amsmath,amssymb}
\usepackage{booktabs,longtable,tabularx,array}
\usepackage{enumitem}
\usepackage{listings}
\usepackage[most]{tcolorbox}
\usepackage{xcolor}
\usepackage{fancyhdr}
\usepackage{titlesec}
\usepackage{needspace}
\usepackage[hidelinks]{hyperref}

\setmainfont{Times New Roman}
\setsansfont{Arial}
\setmonofont{Menlo}[Scale=0.82]
\setCJKmainfont{Songti TC}
\setCJKsansfont{Heiti TC}
\setCJKmonofont{Heiti TC}

\definecolor{Navy}{HTML}{17365D}
\definecolor{Blue}{HTML}{2563A6}
\definecolor{Teal}{HTML}{087F8C}
\definecolor{Green}{HTML}{2E7D4F}
\definecolor{Orange}{HTML}{B45309}
\definecolor{Red}{HTML}{A63D40}
\definecolor{Ink}{HTML}{202733}
\definecolor{Muted}{HTML}{5F6B7A}
\definecolor{Paper}{HTML}{F6F8FB}
\definecolor{BluePale}{HTML}{EAF2FB}
\definecolor{GreenPale}{HTML}{EAF6EF}
\definecolor{OrangePale}{HTML}{FFF4E5}
\definecolor{RedPale}{HTML}{FCEEEF}
\definecolor{Rule}{HTML}{C7D0DC}
\definecolor{CodeBg}{HTML}{F4F6F8}

\hypersetup{
  colorlinks=true,
  linkcolor=Navy,
  urlcolor=Blue,
  pdftitle={C++ Struct 與 Class 完整參考筆記},
  pdfauthor={Week 2 程式設計教材}
}

\setlength{\parindent}{0pt}
\setlength{\parskip}{0.48em}
\linespread{1.14}
\setlist[itemize]{leftmargin=1.7em,itemsep=0.25em,topsep=0.3em}
\setlist[enumerate]{leftmargin=1.9em,itemsep=0.3em,topsep=0.3em}
\renewcommand{\arraystretch}{1.35}
\setcounter{tocdepth}{2}
\renewcommand{\contentsname}{目錄}

\pagestyle{fancy}
\fancyhf{}
\fancyhead[L]{\small\color{Muted}程式設計｜Week 2}
\fancyhead[R]{\small\color{Muted}C++ \texttt{struct} 與 \texttt{class}}
\fancyfoot[C]{\small\color{Muted}\thepage}
\renewcommand{\headrulewidth}{0.3pt}
\renewcommand{\headrule}{\hbox to\headwidth{\color{Rule}\leaders\hrule height \headrulewidth\hfill}}

\titleformat{\section}{\Large\bfseries\sffamily\color{Navy}}{\thesection}{0.65em}{}
\titleformat{\subsection}{\large\bfseries\sffamily\color{Blue}}{\thesubsection}{0.6em}{}
\titleformat{\subsubsection}{\normalsize\bfseries\sffamily\color{Ink}}{\thesubsubsection}{0.55em}{}
\titlespacing*{\section}{0pt}{1.8em}{0.7em}
\titlespacing*{\subsection}{0pt}{1.3em}{0.45em}

\lstdefinestyle{cpp}{
  language=C++,
  basicstyle=\ttfamily\fontsize{9.3}{12.2}\selectfont,
  keywordstyle=\bfseries\color{Navy},
  commentstyle=\itshape\color{Green},
  stringstyle=\color{Orange},
  numbers=left,
  numberstyle=\tiny\color{Muted},
  numbersep=8pt,
  showstringspaces=false,
  keepspaces=true,
  columns=fullflexible,
  breaklines=true,
  tabsize=4,
  frame=single,
  rulecolor=\color{Rule},
  backgroundcolor=\color{CodeBg},
  framesep=6pt,
  xleftmargin=1.1em,
  xrightmargin=0.3em,
  aboveskip=0.65em,
  belowskip=0.65em
}
\lstset{style=cpp}

\newtcolorbox{keybox}[1]{
  breakable,colback=BluePale,colframe=Blue,boxrule=0.6pt,
  arc=2mm,left=2.5mm,right=2.5mm,top=1.8mm,bottom=1.8mm,
  fonttitle=\bfseries\sffamily,title=#1
}
\newtcolorbox{examplebox}[1]{
  breakable,colback=GreenPale,colframe=Green,boxrule=0.6pt,
  arc=2mm,left=2.5mm,right=2.5mm,top=1.8mm,bottom=1.8mm,
  fonttitle=\bfseries\sffamily,title=#1
}
\newtcolorbox{warnbox}[1]{
  breakable,colback=RedPale,colframe=Red,boxrule=0.6pt,
  arc=2mm,left=2.5mm,right=2.5mm,top=1.8mm,bottom=1.8mm,
  fonttitle=\bfseries\sffamily,title=#1
}
\newtcolorbox{skipbox}[1]{
  breakable,colback=OrangePale,colframe=Orange,boxrule=0.6pt,
  arc=2mm,left=2.5mm,right=2.5mm,top=1.8mm,bottom=1.8mm,
  fonttitle=\bfseries\sffamily,title=#1
}

\newcommand{\term}[2]{\textbf{#1}（\textit{#2}）}
\newcommand{\code}[1]{\texttt{#1}}
\newcommand{\yes}{\textcolor{Green}{\textbf{可以}}}
\newcommand{\no}{\textcolor{Red}{\textbf{不可以}}}

\begin{document}

\begin{titlepage}
  \pagecolor{Paper}
  \vspace*{20mm}
  {\sffamily\bfseries\color{Blue}\Large WEEK 2｜完整參考筆記}\par
  \vspace{10mm}
  {\sffamily\bfseries\color{Navy}\fontsize{30}{37}\selectfont
    C++ \texttt{struct} 與 \texttt{class}}
  \par\vspace{6mm}
  {\sffamily\color{Ink}\Large
    從「把資料放在一起」到「設計安全、好用的自訂型別」}
  \par\vspace{17mm}

  \begin{tcolorbox}[
    colback=white,colframe=Rule,boxrule=0.7pt,arc=3mm,
    left=5mm,right=5mm,top=4mm,bottom=4mm]
    \textbf{讀者設定}\quad 已經會變數、\code{if}、迴圈與函式；第一次系統學習自訂型別。\\[1mm]
    \textbf{語言版本}\quad 主線範例使用 C++17；C++20 專屬語法會特別標示。\\[1mm]
    \textbf{閱讀目標}\quad 看得懂、寫得出、能預測輸出，也能解釋為什麼這樣設計。
  \end{tcolorbox}

  \vfill
  {\color{Muted}\small
  內容順序：重點 → 白話解釋 → 小例子 → 必要細節 → 練習與答案}
\end{titlepage}
\nopagecolor

\tableofcontents
\clearpage

\section*{先從這一頁開始}
\addcontentsline{toc}{section}{先從這一頁開始}

\begin{keybox}{一句話抓住兩個主題}
\code{struct} 與 \code{class} 都是在建立一種新的資料型別。這個型別可以同時描述「物件裡有什麼資料」和「物件可以做什麼事」。兩者最直接的語言差異，是未寫存取權限時，\code{struct} 預設 \code{public}，\code{class} 預設 \code{private}。
\end{keybox}

先看一個最小例子。輸入是學生的姓名與分數；型別把兩項資料綁成一組；輸出是完整學生資料。

\begin{lstlisting}
#include <iostream>
#include <string>
using namespace std;

struct Student {
    string name;
    int score;

    bool passed() const { return score >= 60; }
};

int main() {
    Student s{"Mina", 83};
    cout << s.name << ' ' << s.score << ' '
         << s.passed() << '\n';
}
\end{lstlisting}

執行順序是：
\begin{enumerate}
  \item \code{Student} 宣告一種新型別，每個物件都有 \code{name} 與 \code{score}。
  \item \code{Student s\{"Mina", 83\};} 建立一個物件 \code{s}。
  \item \code{s.name} 用點號讀取這個物件的姓名。
  \item \code{s.passed()} 呼叫 \code{s} 自己的成員函式，檢查 \code{s.score} 是否至少 60。
  \item 輸出是 \code{Mina 83 1}；最後的 1 代表 \code{true}。
\end{enumerate}

\begin{skipbox}{建議閱讀路線}
\textbf{第一次讀：}第 1、2、3、5、6、7、8 節，先略過標示「進階查閱」的段落。\\
\textbf{準備考試：}再讀第 4、9、10、11、12 節與速查表。\\
\textbf{已會基礎：}直接讀第 8 節的建構／解構、第 10 節複製語意、第 11 節運算子多載。
\end{skipbox}

\section{先建立物件與型別的心智模型}

\subsection{型別、物件、成員各自是什麼？}

\begin{longtable}{>{\bfseries}p{28mm}p{47mm}p{73mm}}
\toprule
名稱 & 在本程式中的角色 & 具體例子 \\
\midrule
\endhead
型別 type & 規定資料長什麼樣、允許哪些操作 & \code{Student} \\
物件 object & 依照型別建立的一份實際資料 & \code{Student s\{"Mina", 83\};} 中的 \code{s} \\
資料成員 data member & 每個物件自己保存的欄位 & \code{name}、\code{score} \\
成員函式 member function & 由物件呼叫、可以使用該物件資料的函式 & \code{passed()} \\
介面 interface & 類別外部被允許使用的操作方式 & 公開的建構子、\code{deposit()}、\code{balance()} \\
\bottomrule
\end{longtable}

\begin{examplebox}{小例子：同一型別可以建立很多物件}
\code{Student a\{"Amy", 90\};} 與 \code{Student b\{"Bo", 55\};} 都依照 \code{Student} 建立，但兩者各自保存姓名與分數。修改 \code{a.score} 不會自動改到 \code{b.score}。
\end{examplebox}

\subsection{為什麼不用平行陣列就好？}

若姓名和分數分開存，兩個容器的相同索引必須永遠代表同一人：

\begin{lstlisting}
vector<string> names = {"Amy", "Bo"};
vector<int> scores = {90, 55};
\end{lstlisting}

一旦只刪到其中一邊，資料就錯配。改成 \code{vector<Student>} 後，每個學生的姓名與分數會一起移動：

\begin{lstlisting}
vector<Student> students = {{"Amy", 90}, {"Bo", 55}};
\end{lstlisting}

\section{struct：把相關資料組成一個型別}

\subsection{宣告、建立與點號存取}

\begin{lstlisting}
struct Point {
    int x;
    int y;
};  // 型別宣告結尾需要分號

Point p{3, 4};
cout << p.x << ' ' << p.y;  // 3 4
p.x = 10;                   // p 變成 (10, 4)
\end{lstlisting}

輸入 \code{\{3, 4\}} 依欄位宣告順序放入 \code{x}、\code{y}。\code{p.x} 中的點號表示「從物件 \code{p} 取得成員 \code{x}」。

\begin{warnbox}{最常見的第一個錯誤}
\code{struct Point \{ ... \}} 後面少了分號，編譯器常在下一行才報錯。看到「expected ';'」時，要回頭檢查型別定義結尾。
\end{warnbox}

\subsection{初始化：有給值、預設值、沒給值}

\begin{lstlisting}
struct Product {
    string name = "未命名";  // 預設成員初始值
    int price = 0;
};

Product a;                   // name="未命名", price=0
Product b{"筆", 25};        // 指定兩個欄位
Product c{};                 // 使用預設成員初始值
\end{lstlisting}

若內建型別欄位沒有預設值，\code{Product a;} 之類的寫法可能留下未初始化數值。對初學者最安全的習慣是：提供合理的預設成員初始值，或用大括號確實給值。

\begin{skipbox}{C++20 補充：指定欄位名稱}
C++20 可以寫 \code{Point p\{.x = 3, .y = 4\};}。本講義主線使用 C++17，因此範例不依賴這種 designated initializer（指定初始化）語法。
\end{skipbox}

\subsection{巢狀 struct、陣列與 vector}

\begin{lstlisting}
struct Date {
    int year;
    int month;
    int day;
};

struct Event {
    string title;
    Date date;  // 一個 Event 內含一個 Date
};

Event exam{"期中考", {2026, 11, 5}};
cout << exam.date.month;  // 11

Student room[2] = {{"Amy", 90}, {"Bo", 55}};
vector<Student> group = {{"Cy", 72}, {"Di", 88}};
\end{lstlisting}

\code{exam.date.month} 是連續兩次成員存取：先從 \code{exam} 取得 \code{date}，再從那個 \code{Date} 取得 \code{month}。

\subsection{指標使用箭頭運算子}

\begin{lstlisting}
Student s{"Mina", 83};
Student* ptr = &s;

cout << ptr->name;     // 等同 (*ptr).name
ptr->score = 90;       // 修改的是原本的 s.score
\end{lstlisting}

\code{ptr} 存的是 \code{s} 的位址。左邊是指標時使用 \code{->}。例如 \code{ptr->score} 等同 \code{(*ptr).score}。

\section{struct 與函式如何配合}

\subsection{傳值、傳參考、傳 const 參考}

\begin{longtable}{>{\ttfamily}p{45mm}p{42mm}p{58mm}}
\toprule
參數寫法 & 是否複製 & 函式能否改呼叫端原物件 \\
\midrule
\endhead
Student s & 會建立一份副本 & 不會；只改到副本 \\
Student\& s & 不複製 & 會 \\
const Student\& s & 不複製 & 不會；適合只讀 \\
\bottomrule
\end{longtable}

\begin{lstlisting}
void addBonus(Student s) {       // 傳值：只改副本
    s.score += 5;
}

void addBonusInPlace(Student& s) { // 傳參考：改原物件
    s.score += 5;
}

void printStudent(const Student& s) { // 不複製，只讀
    cout << s.name << ' ' << s.score;
}
\end{lstlisting}

若 \code{Student amy\{"Amy", 80\};}：
\begin{itemize}
  \item 呼叫 \code{addBonus(amy)} 後，\code{amy.score} 仍是 80。
  \item 呼叫 \code{addBonusInPlace(amy)} 後，\code{amy.score} 變成 85。
  \item \code{printStudent(amy)} 不複製，也不允許函式修改 \code{amy}。
\end{itemize}

\subsection{自訂型別也可以當回傳值}

\begin{lstlisting}
Student better(const Student& a, const Student& b) {
    if (a.score >= b.score) return a;
    return b;
}
\end{lstlisting}

輸入是兩個唯讀的 \code{Student}；函式比較分數；輸出是一個新的 \code{Student} 值。同分時因為條件是 \code{>=}，會回傳第一位。

\subsection{成員函式：讓操作靠近資料}

\begin{lstlisting}
struct Rectangle {
    double width = 0;
    double height = 0;

    double area() const {
        return width * height;
    }
};

Rectangle r{3.0, 4.0};
cout << r.area();  // 12
\end{lstlisting}

\code{area()} 不需要另外接收 \code{Rectangle} 參數，因為它由 \code{r} 呼叫，能直接使用 \code{r.width} 與 \code{r.height}。函式尾端的 \code{const} 表示它承諾不修改呼叫它的物件。

\section{struct 和 class 到底差在哪裡？}

\begin{keybox}{已由 C++ 語言直接規定的差異}
\begin{itemize}
  \item \code{struct} 的成員若未標示權限，預設是 \code{public}。
  \item \code{class} 的成員若未標示權限，預設是 \code{private}。
  \item 未寫繼承方式時，\code{struct} 預設公開繼承，\code{class} 預設私有繼承。
\end{itemize}
\end{keybox}

兩者都可以有：資料成員、成員函式、建構子、解構子、\code{public/private/protected}、運算子多載、繼承、樣板。所以下列說法是錯的：「\code{struct} 只能放資料，\code{class} 才能放函式。」

\begin{longtable}{>{\bfseries}p{31mm}p{56mm}p{56mm}}
\toprule
比較項目 & 常見 struct 用法 & 常見 class 用法 \\
\midrule
\endhead
設計意圖 & 簡單資料組合，欄位可直接讀寫 & 維持規則，外部透過介面操作 \\
預設成員權限 & public & private \\
是否能有方法 & 可以 & 可以 \\
是否能有建構子 & 可以 & 可以 \\
編譯器本質 & 都能定義自訂型別 & 都能定義自訂型別 \\
\bottomrule
\end{longtable}

\begin{examplebox}{設計判斷}
\code{Point\{x,y\}} 若只是交換資料，可以用 \code{struct}。銀行帳戶必須拒絕負數存款、拒絕超額提款，就適合用 \code{class} 隱藏餘額，再透過方法檢查規則。這是設計慣例，不是硬性規定。
\end{examplebox}

\section{class 與封裝：把規則守在型別裡}

\subsection{public、private、protected}

\begin{lstlisting}
class BankAccount {
private:
    string owner;
    int balanceInDollars = 0;

public:
    bool deposit(int amount);
    bool withdraw(int amount);
    int balance() const;
};
\end{lstlisting}

\begin{itemize}
  \item \code{public}：類別外可以使用，是提供給使用者的介面。
  \item \code{private}：只有類別自己的成員函式與被允許的 friend 能直接存取。
  \item \code{protected}：類別自己與衍生類別可直接存取；初學階段通常先少用。
\end{itemize}

\term{封裝}{encapsulation} 是把資料和操作資料的函式放在同一型別，並限制外部只能透過公開介面改變狀態。它的直接效果是：輸入檢查只需集中在少數方法中。

\begin{lstlisting}
bool BankAccount::deposit(int amount) {
    if (amount <= 0) return false;
    balanceInDollars += amount;
    return true;
}

bool BankAccount::withdraw(int amount) {
    if (amount <= 0 || amount > balanceInDollars) return false;
    balanceInDollars -= amount;
    return true;
}

int BankAccount::balance() const {
    return balanceInDollars;
}
\end{lstlisting}

若餘額原本是 100：\code{withdraw(40)} 回傳 \code{true} 且餘額變 60；\code{withdraw(80)} 回傳 \code{false} 且餘額維持 60。

\subsection{類別內宣告、類別外定義與 ::}

\code{bool deposit(int amount);} 放在類別內，先告訴編譯器函式的名稱、輸入與輸出。實作前面的 \code{BankAccount::} 是 scope resolution operator（範圍解析運算子），表示這個 \code{deposit} 屬於 \code{BankAccount}。

\begin{lstlisting}
// BankAccount.h：通常放宣告
class BankAccount {
public:
    bool deposit(int amount);
};

// BankAccount.cpp：通常放定義
bool BankAccount::deposit(int amount) {
    // ...
}
\end{lstlisting}

小型練習可以全部放在一個 \code{.cpp}。專案變大後，分開宣告與定義可以縮短閱讀範圍，也讓多個來源檔共用同一介面。

\subsection{getter 與 setter 不是越多越好}

getter 是讀取方法，setter 是修改方法。若把每個 private 欄位都做成無條件 getter/setter，外部仍可能任意改壞狀態。

\begin{lstlisting}
bool setScore(int newScore) {
    if (newScore < 0 || newScore > 100) return false;
    score = newScore;
    return true;
}
\end{lstlisting}

這個 setter 的輸入是新分數；先檢查 0 到 100；合法才改欄位並回傳 \code{true}；非法時回傳 \code{false}，而且保留原值。

\section{成員函式、const 與 this}

\subsection{函式後面的 const 保護誰？}

\begin{lstlisting}
class Counter {
private:
    int value_ = 0;

public:
    void add(int amount) { value_ += amount; }
    int value() const { return value_; }
};
\end{lstlisting}

\code{add} 會修改呼叫者，所以不能加尾端 \code{const}。\code{value()} 只讀取資料，所以加上尾端 \code{const}。這讓 \code{const Counter fixed;} 可以呼叫 \code{fixed.value()}，卻不能呼叫 \code{fixed.add(1)}。

\begin{warnbox}{三個 const 的位置不要混在一起}
\code{const Student\& s}：參數 \code{s} 不能被修改。\\
\code{int score() const}：成員函式不能修改呼叫它的物件。\\
\code{const int* p}：不能透過指標 \code{p} 修改它所指的整數。
\end{warnbox}

\subsection{const 多載：一般物件可改，const 物件只讀}

\begin{lstlisting}
class Pair {
private:
    int data_[2] = {0, 0};

public:
    int& at(int index) { return data_[index]; }
    const int& at(int index) const { return data_[index]; }
};
\end{lstlisting}

一般物件呼叫第一個版本，所以 \code{p.at(0) = 9;} 能修改欄位。\code{const Pair} 只能呼叫第二個版本，回傳唯讀參考，因此不能放在指定運算的左邊。

\subsection{this 指向目前的呼叫者}

\begin{lstlisting}
class Counter {
private:
    int value_ = 0;

public:
    Counter& add(int amount) {
        this->value_ += amount;
        return *this;
    }
};
\end{lstlisting}

若執行 \code{a.add(2)}，\code{this} 指向 \code{a}，\code{*this} 就是 \code{a} 本身。回傳 \code{Counter\&} 讓 \code{a.add(2).add(3)} 連續操作同一物件。平常寫 \code{value\_} 就足夠；參數和欄位同名時，\code{this->value\_} 才特別有助於區分。

\section{建構子與解構子：物件怎麼開始、怎麼結束}

\subsection{預設、參數化與多載建構子}

\begin{lstlisting}
class Student {
private:
    string name_;
    int score_ = 0;

public:
    Student() : name_("未命名") {}

    Student(string name) : name_(name) {}

    Student(string name, int score)
        : name_(name) {
        setScore(score);
    }

    bool setScore(int score) {
        if (score < 0 || score > 100) return false;
        score_ = score;
        return true;
    }
};
\end{lstlisting}

同名但參數不同的建構子是 overloading（多載）。編譯器會依呼叫時提供的參數選版本：

\begin{lstlisting}
Student a;               // 預設建構子
Student b("Amy");       // 一個 string
Student c("Bo", 85);    // string 加 int
\end{lstlisting}

\subsection{成員初始化列表}

\begin{lstlisting}
class Point {
    const int id_;
    int x_;
    int y_;

public:
    Point(int id, int x, int y)
        : id_(id), x_(x), y_(y) {}
};
\end{lstlisting}

冒號後面是 member initializer list（成員初始化列表）。它在成員建立時直接給初值；建構子主體中的指定，是成員建立後才重新指定。\code{const} 成員、參考成員，以及沒有預設建構子的物件成員，都必須透過初始化列表處理。

\begin{warnbox}{初始化順序看宣告，不看列表順序}
成員實際初始化順序是它們在類別內的宣告順序。即使寫 \code{: y\_(...), x\_(...)}，仍會先初始化先宣告的 \code{x\_}。為避免誤會，初始化列表最好也照宣告順序寫。
\end{warnbox}

\subsection{delegating constructor：讓建構子重用建構子}

\begin{lstlisting}
class Point {
    int x_;
    int y_;

public:
    Point() : Point(0, 0) {}
    Point(int x, int y) : x_(x), y_(y) {}
};
\end{lstlisting}

\code{Point()} 把工作交給 \code{Point(int,int)}。這比兩個建構子各寫一次初始化規則更不容易不一致。

\subsection{explicit：避免單一參數被偷偷轉型}

\begin{lstlisting}
class Meter {
    double value_;
public:
    explicit Meter(double value) : value_(value) {}
};

Meter a(3.5);       // 可以
// Meter b = 3.5;   // explicit 會阻止隱式轉換
\end{lstlisting}

單一參數建構子可能被編譯器當成「從參數型別轉成此類別」的方法。若這不是你要的行為，就加 \code{explicit}，要求呼叫端清楚寫出型別。

\subsection{解構子與 RAII}

\begin{lstlisting}
class Trace {
    string name_;
public:
    explicit Trace(string name) : name_(name) {
        cout << "建立 " << name_ << '\n';
    }

    ~Trace() {
        cout << "離開 " << name_ << '\n';
    }
};
\end{lstlisting}

解構子與類別同名、前面有 \code{\~}、沒有參數，也沒有回傳型別。物件離開生命週期時自動執行。RAII 是「資源取得即初始化」：讓資源由物件保存，建構時取得，解構時自動釋放。實務上優先使用 \code{string}、\code{vector}、智慧指標等已經管理資源的型別。

\begin{warnbox}{不要手動呼叫一般區域物件的解構子}
區域物件離開區塊時會自動解構。手動呼叫後，系統之後又自動呼叫，可能重複釋放資源。
\end{warnbox}

\section{static、friend、巢狀型別與組合}

\subsection{static 成員屬於整個類別}

\begin{lstlisting}
class Ticket {
private:
    int id_;
    static int nextId_;

public:
    Ticket() : id_(nextId_++) {}
    int id() const { return id_; }
    static int nextId() { return nextId_; }
};

int Ticket::nextId_ = 1;  // C++17 也可改用 inline static
\end{lstlisting}

每個物件各有自己的 \code{id\_}，但所有 \code{Ticket} 共用一份 \code{nextId\_}。static 成員函式沒有特定呼叫者，因此沒有 \code{this}，也不能直接讀取某一個物件的非 static 欄位。

\subsection{friend 是精準授權，不是新的權限層級}

\begin{lstlisting}
class Point {
    int x_;
    int y_;
public:
    Point(int x, int y) : x_(x), y_(y) {}
    friend ostream& operator<<(ostream& out, const Point& p);
};

ostream& operator<<(ostream& out, const Point& p) {
    return out << '(' << p.x_ << ", " << p.y_ << ')';
}
\end{lstlisting}

\code{operator<<} 不是成員函式，但因為被宣告成 friend，可以讀取 \code{Point} 的 private 成員。friend 會增加耦合，只有外部函式確實需要看內部資料時才使用。

\subsection{組合：一個物件擁有另一個物件}

\begin{lstlisting}
class Engine {
public:
    void start() { cout << "engine on\n"; }
};

class Car {
private:
    Engine engine_;  // Car 擁有一個 Engine
public:
    void start() { engine_.start(); }
};
\end{lstlisting}

composition（組合）表示「has-a」：\code{Car} 有一個 \code{Engine}。\code{Car} 建立時，\code{engine\_} 先建立；\code{Car} 解構時，成員會自動解構。能用組合表達時，通常比為了重用幾行程式而繼承更清楚。

\section{複製、指定、移動與資源管理}

\subsection{先分清楚三個看似相近的動作}

\begin{lstlisting}
Point a(1, 2);
Point b = a;   // 建立 b：copy constructor（複製建構子）
Point c(0, 0);
c = a;         // c 已存在：copy assignment（複製指定）
\end{lstlisting}

對只有 \code{int}、\code{string}、\code{vector} 等正常成員的類別，編譯器產生的複製通常已經正確。這叫 Rule of Zero：若成員自己會管理資源，就讓編譯器合成特殊成員函式，不必手寫。

\subsection{淺層複製為何可能出錯？}

\begin{lstlisting}
class BadBuffer {
    int* data_;
public:
    explicit BadBuffer(int value) : data_(new int(value)) {}
    ~BadBuffer() { delete data_; }
};
\end{lstlisting}

若直接 \code{BadBuffer b = a;}，編譯器預設只複製指標位址，\code{a.data\_} 與 \code{b.data\_} 指向同一格記憶體。兩個物件解構時會對同一位址 \code{delete} 兩次，行為未定義。

\begin{keybox}{Rule of Three、Five、Zero}
\begin{itemize}
  \item 若類別必須自己寫解構子、複製建構子、複製指定，通常三者都要一起思考：Rule of Three。
  \item C++11 後再加移動建構子與移動指定：Rule of Five。
  \item 最推薦的方向是讓 \code{vector}、\code{string}、\code{unique\_ptr} 等成員管理資源，讓類別不必手寫這五個函式：Rule of Zero。
\end{itemize}
\end{keybox}

\subsection{進階查閱：完整的深層複製骨架}

\begin{lstlisting}
class Buffer {
    size_t size_ = 0;
    int* data_ = nullptr;

public:
    explicit Buffer(size_t size)
        : size_(size), data_(new int[size]{}) {}

    ~Buffer() { delete[] data_; }

    Buffer(const Buffer& other)
        : size_(other.size_), data_(new int[other.size_]) {
        copy(other.data_, other.data_ + size_, data_);
    }

    Buffer& operator=(const Buffer& other) {
        if (this == &other) return *this;
        Buffer temp(other);
        swap(size_, temp.size_);
        swap(data_, temp.data_);
        return *this;
    }
};
\end{lstlisting}

複製建構子配置一塊新陣列，再複製每個元素，所以兩個物件不共用同一塊陣列。指定運算先檢查自我指定，再用 copy-and-swap（先做安全副本，再交換）維持例外安全。真正專案通常直接使用 \code{vector<int>}，把這些工作交給標準函式庫。

\subsection{move 的目的}

move semantics（移動語意）讓新物件接手暫時物件的資源，不必逐項複製大量資料。參數型別常寫成 \code{T\&\&}。初學者先掌握兩點即可：
\begin{itemize}
  \item 標準容器已正確實作 move；使用 \code{vector} 通常不必自己處理。
  \item 被移動後的物件仍可安全解構、可以再指定，但原內容通常不能再假設。
\end{itemize}

\section{運算子多載：讓自訂型別有一致的操作方式}

\subsection{能做什麼、不能做什麼}

operator overloading（運算子多載）能替自訂型別定義現有運算子的行為，例如 \code{==}、\code{+}、\code{[]}、\code{<<}。它不能創造新符號、不能改變優先順序、不能改變運算元個數，也至少要有一個運算元是自訂型別。

\subsection{== 與 !=：兩邊都不應被修改}

\begin{lstlisting}
bool operator==(const Point& other) const {
    return x_ == other.x_ && y_ == other.y_;
}

bool operator!=(const Point& other) const {
    return !(*this == other);
}
\end{lstlisting}

\code{a == b} 可看成 \code{a.operator==(b)}。參數的 \code{const\&} 保護右邊並避免複製；尾端 \code{const} 保護左邊。

\subsection{+= 修改左邊，+ 產生新結果}

\begin{lstlisting}
Point& operator+=(const Point& other) {
    x_ += other.x_;
    y_ += other.y_;
    return *this;
}

Point operator+(const Point& other) const {
    Point result(*this);
    result += other;
    return result;
}
\end{lstlisting}

\code{a += b} 會改 \code{a}，所以回傳 \code{Point\&} 指向原本左邊。\code{a + b} 不改兩邊，先複製左邊到 \code{result}，再重用 \code{+=}，最後回傳新的值。不能回傳 \code{Point\&} 指向區域變數 \code{result}，因為函式結束後它就不存在。

\subsection{[]：回傳參考才可以放在等號左邊}

\begin{lstlisting}
int& operator[](size_t index) {
    if (index == 0) return x_;
    if (index == 1) return y_;
    throw out_of_range("Point index must be 0 or 1");
}

const int& operator[](size_t index) const {
    if (index == 0) return x_;
    if (index == 1) return y_;
    throw out_of_range("Point index must be 0 or 1");
}
\end{lstlisting}

第一個版本讓 \code{p[0] = 9;} 修改 \code{x\_}。第二個版本讓 \code{const Point} 可以讀，卻不能寫。非法索引丟出例外，而不是讀取不存在的記憶體。

\subsection{<< 為何通常寫在類別外？}

\begin{lstlisting}
ostream& operator<<(ostream& out, const Point& p) {
    out << '(' << p.x() << ", " << p.y() << ')';
    return out;
}
\end{lstlisting}

\code{cout << p} 的左邊是 \code{ostream}，不是 \code{Point}，所以不能做成 \code{Point} 的成員函式。回傳 \code{ostream\&} 才能繼續串接下一個輸出。如果公開 getter 足夠，這個函式不必是 friend。

\subsection{運算子設計檢查表}

\begin{longtable}{p{20mm}p{40mm}p{42mm}p{40mm}}
\toprule
運算子 & 是否改左邊 & 常見回傳型別 & 常見 const 設計 \\
\midrule
\endhead
\code{==, !=} & 否 & \code{bool} & 參數 \code{const\&}，方法 \code{const} \\
\code{+} & 否 & \code{T} & 參數 \code{const\&}，方法 \code{const} \\
\code{+=} & 是 & \code{T\&} & 參數 \code{const\&}，方法不可 \code{const} \\
\code{[]} 讀取 & 否 & \code{const T\&} 或 \code{T} & 方法 \code{const} \\
\code{[]} 寫入 & 可能 & \code{T\&} & 方法不可 \code{const} \\
\code{<<} & 否 & \code{ostream\&} & 物件參數 \code{const\&} \\
\bottomrule
\end{longtable}

\section{繼承與多型：class 的進階地圖}

\begin{skipbox}{這節的定位}
繼承與執行期多型也是 class 能力，但通常應在學生熟悉封裝、參考與生命週期後再深入。這裡提供完整地圖與安全骨架，不把它混進第一輪主線。
\end{skipbox}

\subsection{public inheritance 表示 is-a}

\begin{lstlisting}
class Shape {
public:
    virtual double area() const = 0;
    virtual ~Shape() = default;
};

class Circle : public Shape {
    double radius_;
public:
    explicit Circle(double radius) : radius_(radius) {}
    double area() const override {
        return 3.141592653589793 * radius_ * radius_;
    }
};
\end{lstlisting}

\code{Circle : public Shape} 表示 Circle 是一種 Shape。\code{virtual} 讓經由基底類別參考或指標呼叫時，執行實際物件的版本。\code{= 0} 是純虛擬函式，讓 \code{Shape} 成為抽象類別，不能直接建立物件。\code{override} 要求編譯器檢查是否真的覆寫成功。

\subsection{多型需要參考或指標}

\begin{lstlisting}
void printArea(const Shape& shape) {
    cout << shape.area() << '\n';
}

Circle c(2.0);
printArea(c);  // 呼叫 Circle::area()
\end{lstlisting}

若用值把衍生物件存成基底物件，衍生部分會被切掉，稱為 object slicing（物件切片）。需要多型時使用基底參考或智慧指標。

\subsection{protected 與繼承權限}

\code{protected} 成員可由衍生類別直接存取，但外部不行。通常仍優先把資料設為 private，透過 protected/public 函式維持規則。public inheritance 會保留基底的 public 介面；private inheritance 比較少用，含義也不是一般的 is-a。

\section{完整整合範例：安全的 BankAccount}

\begin{lstlisting}
#include <iostream>
#include <string>
using namespace std;

class BankAccount {
private:
    string owner_;
    int balance_ = 0;
    static int accountCount_;

public:
    explicit BankAccount(string owner)
        : owner_(owner) {
        ++accountCount_;
    }

    BankAccount(string owner, int initialBalance)
        : BankAccount(owner) {
        deposit(initialBalance);
    }

    ~BankAccount() { --accountCount_; }

    const string& owner() const { return owner_; }
    int balance() const { return balance_; }
    static int accountCount() { return accountCount_; }

    bool deposit(int amount) {
        if (amount <= 0) return false;
        balance_ += amount;
        return true;
    }

    bool withdraw(int amount) {
        if (amount <= 0 || amount > balance_) return false;
        balance_ -= amount;
        return true;
    }

    bool transferTo(BankAccount& target, int amount) {
        if (&target == this) return false;
        if (!withdraw(amount)) return false;
        target.balance_ += amount;
        return true;
    }
};

int BankAccount::accountCount_ = 0;

int main() {
    BankAccount amy("Amy", 100);
    BankAccount bo("Bo", 20);

    cout << amy.transferTo(bo, 40) << '\n';
    cout << amy.balance() << ' ' << bo.balance() << '\n';
    cout << BankAccount::accountCount() << '\n';
}
\end{lstlisting}

預期輸出：
\begin{lstlisting}[language={}]
1
60 60
2
\end{lstlisting}

因果流程：\code{transferTo} 先拒絕自己轉給自己；再呼叫 \code{withdraw} 檢查金額與餘額；提款成功後才增加目標帳戶餘額。成員函式可以讀寫同類別另一個物件的 private 欄位，所以 \code{target.balance\_} 合法。

\section{常見錯誤：看錯誤訊息時先找哪裡？}

\begin{longtable}{p{45mm}p{48mm}p{50mm}}
\toprule
症狀 & 常見原因 & 修正方向 \\
\midrule
\endhead
型別定義下一行報錯 & class/struct 結尾少分號 & 檢查右大括號後是否有 \code{;} \\
private within this context & 類別外直接存取 private & 改用公開方法，或重新設計介面 \\
no matching constructor & 傳入參數與建構子不相符 & 核對數量、型別與 explicit \\
deleted default constructor / no default constructor & 已宣告其他建構子後，編譯器未必提供無參數版本 & 明確寫 \code{T() = default;} 或提供預設建構子 \\
discard qualifiers / const 錯誤 & const 物件呼叫了非 const 方法 & 若方法只讀資料，補尾端 \code{const} \\
multiple definition & 在標頭重複定義非 inline 的 static 資料或函式 & 宣告、定義分開；或 C++17 用 \code{inline static} \\
程式偶爾出現怪值 & 內建型別欄位未初始化 & 給預設成員初始值或建構子初始化 \\
double free / crash & 擁有裸指標卻只做淺層複製 & 優先改用 vector/智慧指標，或正確實作 Rule of Five \\
結果改不到原物件 & 參數用傳值，改到副本 & 若確實要改呼叫端，使用非 const 參考 \\
回傳後資料壞掉 & 回傳區域變數的參考或指標 & 區域結果應回傳值 \\
\bottomrule
\end{longtable}

\section{練習：從會讀到會設計}

\subsection*{練習 1｜struct 與函式}
宣告 \code{Book}，包含書名、頁數與價格。寫 \code{printBook(const Book\&)}，以及 \code{cheaper(const Book\&, const Book\&)}，同價時回傳第一本。

\subsection*{練習 2｜封裝與邊界}
設計 \code{Temperature} 類別，內部保存攝氏溫度。建構子收到低於絕對零度 $-273.15$ 的輸入時，改存 $-273.15$。提供 \code{celsius()} 與 \code{fahrenheit()}；兩個都不可修改物件。

\subsection*{練習 3｜建構子與 setter}
設計 \code{class Student}。分數必須是 0 到 100；非法建構值保留 0；\code{setScore} 失敗時保留原值並回傳 \code{false}。

\subsection*{練習 4｜const 與參考}
解釋下列三個函式的差異，並寫出每個函式是否會複製、是否能改呼叫端：
\begin{lstlisting}
void f(Student s);
void g(Student& s);
void h(const Student& s);
\end{lstlisting}

\subsection*{練習 5｜運算子多載}
替 \code{Vec2} 實作 \code{==}、\code{+=}、\code{+}、兩個版本的 \code{[]} 與輸出 \code{<<}。測試 \code{a += a}、\code{a + b + c}、\code{const Vec2} 讀索引、非法索引。

\subsection*{練習 6｜獨立設計題}
設計 \code{class InventoryItem}：名稱、單價、庫存量是 private。補貨數量必須為正；出售數量必須為正且不可超過庫存；總價以「單價乘庫存」計算。列出至少 6 個邊界測試。

\clearpage
\section{練習參考答案}

\subsection*{答案 1}
\begin{lstlisting}
struct Book {
    string title;
    int pages;
    int price;
};

void printBook(const Book& book) {
    cout << book.title << ' ' << book.pages << ' ' << book.price;
}

Book cheaper(const Book& a, const Book& b) {
    return (a.price <= b.price) ? a : b;
}
\end{lstlisting}

\subsection*{答案 2}
\begin{lstlisting}
class Temperature {
    double celsius_;
public:
    explicit Temperature(double celsius)
        : celsius_(celsius < -273.15 ? -273.15 : celsius) {}

    double celsius() const { return celsius_; }
    double fahrenheit() const { return celsius_ * 9.0 / 5.0 + 32.0; }
};
\end{lstlisting}

\subsection*{答案 3}
\begin{lstlisting}
class Student {
    string name_;
    int score_ = 0;
public:
    Student(string name, int score) : name_(name) {
        setScore(score);
    }

    bool setScore(int score) {
        if (score < 0 || score > 100) return false;
        score_ = score;
        return true;
    }

    int score() const { return score_; }
};
\end{lstlisting}

\subsection*{答案 4}
\code{f} 傳值，會複製，不能改呼叫端原物件；\code{g} 傳非 const 參考，不複製而且可以改原物件；\code{h} 傳 const 參考，不複製也不能改原物件。只讀取大型物件時，\code{h} 通常最合適。

\subsection*{答案 5：核心類別}
\begin{lstlisting}
class Vec2 {
    int x_;
    int y_;
public:
    Vec2(int x, int y) : x_(x), y_(y) {}

    int x() const { return x_; }
    int y() const { return y_; }

    bool operator==(const Vec2& other) const {
        return x_ == other.x_ && y_ == other.y_;
    }

    Vec2& operator+=(const Vec2& other) {
        x_ += other.x_;
        y_ += other.y_;
        return *this;
    }

    Vec2 operator+(const Vec2& other) const {
        Vec2 result(*this);
        result += other;
        return result;
    }

    int& operator[](size_t index) {
        if (index == 0) return x_;
        if (index == 1) return y_;
        throw out_of_range("Vec2 index must be 0 or 1");
    }

    const int& operator[](size_t index) const {
        if (index == 0) return x_;
        if (index == 1) return y_;
        throw out_of_range("Vec2 index must be 0 or 1");
    }
};

ostream& operator<<(ostream& out, const Vec2& v) {
    return out << '(' << v.x() << ", " << v.y() << ')';
}
\end{lstlisting}

\subsection*{答案 6：測試重點}
至少測：正常補貨、補貨 0、負數補貨、正常出售、出售 0、負數出售、出售剛好等於庫存、出售超過庫存、庫存變成 0 後再出售。每個失敗操作都要確認庫存沒有被改動。

\section{一頁速查表}

\subsection*{基本骨架}
\begin{lstlisting}
struct Data {
    int value = 0;       // 預設 public
};

class SafeType {
private:
    int value_ = 0;
public:
    explicit SafeType(int value) : value_(value) {}
    int value() const { return value_; }
    bool setValue(int value);
};
\end{lstlisting}

\subsection*{選參數方式}
\begin{itemize}
  \item 小型內建值且不需改原值：傳值，例如 \code{int amount}。
  \item 要修改呼叫端：\code{T\&}。
  \item 只讀且避免複製：\code{const T\&}。
\end{itemize}

\subsection*{選回傳方式}
\begin{itemize}
  \item 產生新結果：回傳 \code{T}。
  \item 回傳已存在且生命週期足夠長的物件：才考慮 \code{T\&} 或 \code{const T\&}。
  \item 永遠不要回傳區域變數的參考或指標。
\end{itemize}

\subsection*{設計 class 前的五個問題}
\begin{enumerate}
  \item 每個物件要保存哪些資料？初值是什麼？
  \item 哪些狀態不合法？誰負責阻止它？
  \item 外部真正需要哪些 public 操作？
  \item 哪些方法只讀，應該標成 \code{const}？
  \item 類別是否直接擁有資源？能否改用標準容器以遵守 Rule of Zero？
\end{enumerate}

\section{術語中英對照}

\begin{longtable}{p{43mm}p{48mm}p{48mm}}
\toprule
中文 & 英文 & 在本筆記中的意思 \\
\midrule
\endhead
自訂型別 & user-defined type & 程式設計者用 struct/class 建立的型別 \\
物件 & object & 依型別建立的一份實際資料 \\
實例 & instance & 物件的另一種常見稱呼 \\
資料成員 & data member & 物件保存的欄位 \\
成員函式 & member function & 由物件呼叫、可存取其成員的函式 \\
存取權限 & access control & public/private/protected 的限制 \\
封裝 & encapsulation & 把資料與操作放在一起並限制外部存取 \\
介面 & interface & 類別外可用的公開操作 \\
建構子 & constructor & 物件建立時執行的特殊函式 \\
解構子 & destructor & 物件生命週期結束時執行的特殊函式 \\
成員初始化列表 & member initializer list & 在建構子主體前直接初始化成員 \\
函式多載 & function overloading & 同名、參數不同的多個函式 \\
運算子多載 & operator overloading & 定義運算子用於自訂型別時的行為 \\
組合 & composition & 一個物件擁有另一個物件 \\
繼承 & inheritance & 衍生類別沿用並擴充基底類別介面 \\
多型 & polymorphism & 透過共同介面呼叫不同實際型別的行為 \\
淺層複製 & shallow copy & 只複製位址等表面成員，可能共用資源 \\
深層複製 & deep copy & 建立獨立資源並複製內容 \\
移動語意 & move semantics & 把資源所有權轉交給新物件 \\
\bottomrule
\end{longtable}

\begin{keybox}{最後自我檢查}
確認你能回答：\code{struct}/\code{class} 的預設權限有何差別？\code{private}、初始化列表與尾端 \code{const} 各解決什麼問題？何時使用 \code{T}、\code{T\&}、\code{const T\&}？為什麼 \code{+} 回傳值、\code{+=} 回傳參考？裸指標的預設複製有何風險？
\end{keybox}

\end{document}
